% !TEX root = ../../mixed-sklyanin-cubics.tex
\section{有限台加群の perfect 性と残余圏の非滑らかさ}
\label{mix:sec:perfect-modules}

\subsection{有限台の chart 加群と線束 perfect 性}

以下でいう線束 perfect とは
\[
 \mathcal C_A=\operatorname{thick}\{\mathcal O(j):j\in\mathbb Z\}
 \subset D^b(\operatorname{qgr}A)
\]
に属することをいう。この部分圏の proper 性は、全 coherent 圏の
smoothness を仮定せずに成立する。

\begin{lemma}[有限中心上の有限台 chart 加群]
\label{cmo:lem:finite-support-perfect}
$A$ を次数一で生成される連結局所有限 noetherian 次数付き代数とし、
$u\in Z(A)_1$ を正則元とする。
$u$ を含む可換有限型次数付き中心部分環 $C\subset A$ があり、
$A$ は $C$ 上有限であると仮定する。
$D=(A[u^{-1}])_0$ とし、有限次元右 $D$-加群 $M$ に対して
補題\ref{cmo:lem:chart-ext}の $\mathcal P(M)=\pi(M[u])$ を取る。
このとき
\begin{equation}\label{cmo:eq:chart-perfect-iff}
 \mathcal P(M)\in\mathcal C_A
 \quad\Longleftrightarrow\quad
 \operatorname{pd}_{D}M<\infty.
\end{equation}
\end{lemma}
\begin{proof}
$\ell>0$ を十分割り切れる整数として、$C^{(\ell)}$ が次数一で生成されるようにする。
$Z=\operatorname{Proj}C^{(\ell)}$、
$\mathcal A=\widetilde{A^{(\ell)}}$ と置く。
定理\ref{cmo:thm:relative-smoothness}の証明の有限中心模型と同じく、
\[
 \operatorname{QGr}A\simeq\operatorname{Qcoh}(Z,\mathcal A),\qquad
 \operatorname{qgr}A\simeq\operatorname{coh}(Z,\mathcal A)
\]
がある。これは smoothness の結論に先立つ同値である。
また、次数 $\ell$ の中央元 $f$ を反転した $A[f^{-1}]$ は強次数付きである。
実際、$f$ を次数一の元の長さ $\ell$ の積の和として書き、
先頭または末尾の元を分ければ
$1\in A[f^{-1}]_1A[f^{-1}]_{-1}$ および
$1\in A[f^{-1}]_{-1}A[f^{-1}]_1$ を得る。
従って全ての $\mathcal O(j)$ は局所有限生成射影 $\mathcal A$-加群である。
中央の ample twist $\mathcal A\otimes\mathcal O_Z(m)$ は
$\mathcal O(\ell m)$ に対応し、これらは
$D_{\mathrm{qc}}(Z,\mathcal A)$ を compact に生成する。
有限アファイン被覆上での局所有限射影分解と同じ生成論から、
\begin{equation}\label{cmo:eq:line-is-local-perfect}
 \mathcal C_A=\operatorname{Perf}(Z,\mathcal A)
\end{equation}
である。混合族では、命題\ref{cmo:prop:proper-contraction}の有限線束和 $E$
がこれら全 twist を厚生成するため、同じ等式はその compact generator
の記述からも従う。

ここで $P=M[u]$ 上の中心作用を調べる。
各 $c\in C_e$ は、$M$ 上の $cu^{-e}$ の作用行列に $u^e$ を掛けたものとして作用する。
従って中心像は、$u$ を含む次数付き可換部分環
\[
 S=C/\operatorname{Ann}_C P
 \ \subset\ \operatorname{End}_{\mathbb C}(M)[u]
\]
である。右辺は $\mathbb C[u]$ 上有限であり、
$S$ はその $\mathbb C[u]$-部分加群だから、$S$ も有限である。
従って $S/(u)$ は有限次元であり、
\[
 \operatorname{Proj}S^{(\ell)}\subset Z
\]
は $u\ne0$ 内の有限閉部分スキームである。
より具体的には、$S[u^{-1}]_0\subset\operatorname{End}_{\mathbb C}(M)$ は
有限次元可換代数で、上の Proj はその Spec に等しい。
$\mathcal P(M)$ の中央台はこの有限閉集合に含まれる。

$U=D_+(u^\ell)$ では $\mathcal A|_U$ のアファイン代数は $D$ であり、
$\mathcal P(M)|_U$ は $M$ に対応する。
もし $\operatorname{pd}_D M<\infty$ なら $U$ 上で perfect である。
中央台の補集合では零なので、$Z$ 全体で局所 perfect となる。
式\eqref{cmo:eq:line-is-local-perfect}から順方向を得る。

逆に $j^*:\operatorname{QGr}A\to\operatorname{Mod}D$ は exact で、
$j^*\mathcal O(j)\simeq D$、$j^*\mathcal P(M)=M$ である。
従って $\mathcal P(M)\in\mathcal C_A$ なら
$M\in\operatorname{Perf}(D)$、すなわち有限生成加群 $M$ の射影次元は有限である。
これは、任意の perfect 加群を開埋込みの $j_*$ が perfect に保つという主張ではない。
順方向で不可欠なのは、有限次元性から得た、開 chart 内で閉じた中央台である。
\end{proof}

\begin{proposition}[定常点の線束 perfect 性の完全判定]
\label{cmo:prop:point-perfect}
命題\ref{cmo:prop:character-pd}の $k$ 因子の混合型を取り、
$p_i=(a_i,b_i,c_i)\in(\mathbb C^\times)^3$ とする。
中心商の係数を $c_0=-\sum_i\beta_i\lambda_{p_i}$ とし、
対応する次数付き一点加群 $P=\mathbb C[u]$ と
$\mathcal P=\pi P$ を取る。このとき
\begin{equation}\label{cmo:eq:point-perfect-criterion}
 \mathcal P\in\mathcal C_A
 \quad\Longleftrightarrow\quad
 \text{少なくとも一つの }\kappa(p_i)\ne0.
\end{equation}
有限の場合、対応する右 affine character の射影次元は正確に $3k-1$ である。
特に例\ref{cmo:ex:sixty-four}の全六十四個の定常点は線束 perfect である。
\end{proposition}
\begin{proof}
最初に左右の規約を揃える。各因子に対して、語の順序を逆転させるとともに
\[
 x\longmapsto x,\qquad y\longmapsto z,\qquad
 z\longmapsto y,\qquad u\longmapsto u
\]
とする同型
\[
 \mathcal B_{(a^2,b^2,c^2)}
 \xrightarrow{\sim}
 \mathcal B_{(a^2,c^2,b^2)}^{\mathrm{op}}
\]
がある。三つの関係式を直接代入すれば確認でき、
中心元 $\Omega$ と character はそれぞれ対応する中心元と $(a,c,b)$ へ移る。
$\lambda$ と $\kappa$ は三座標の置換で不変だから、
命題\ref{cmo:prop:character-pd}の左射影次元の結論は右加群にも移る。
有限次元体双対による自己 Ext の比較だけから右射影次元を推論しているのではない。

定理\ref{cmo:thm:finite-center}の普遍中心元の像が生成する中心部分環を使うと、
$A$ は $u$ を含む有限型次数付き中心上有限である。
従って前補題と右 affine character の射影次元判定が適用できる。
一点の場合は中心像がちょうど $\mathbb C[u]$ なので、
中央台は $u\ne0$ 内の一つの閉点である。
六十四点では系\ref{cmo:cor:sixty-four-ext}の二つの不等式により
両因子の $\kappa$ がともに非零である。
\end{proof}

\begin{corollary}[低次元単純加群の coherent な perfect 延長]
\label{cmo:cor:hundred-perfect}
命題\ref{cmo:prop:one-two-dimensional-pd}の百個の左単純加群を体双対して
右加群 $M$ とし、各出現係数の $A_c$ 上で $\mathcal P(M)=\pi(M[u])$ と置く。
これらは全て線束 perfect である。
\end{corollary}
\begin{proof}
体双対は有限次元端点の自己 Ext を、反対向きの自己 Ext に同型に移す。
例えば有限生成射影分解を体双対した単射分解と Hom の随伴からこの同型を得る。
右側の周囲 tensor 代数も有限大域次元六を持ち、$M$ は有限生成なので、
長さ六の有限生成射影分解がある。
従って、双対した後に系\ref{cmo:cor:hypersurface-pd}を右加群へ適用すれば
有限自己 Ext から $\operatorname{pd}M=5$ を得る。
体双対だけから射影次元を移してはいない。
補題\ref{cmo:lem:finite-support-perfect}を適用して結論を得る。
\end{proof}

\begin{remark}[判定の範囲]
この結果は定常点そのものが有限線束複体の直和因子になることを示す。
全 coherent 対象の perfect 性や、全カテゴリー・対角の smoothness を意味しない。
全 $\kappa=0$ の場合には、既存の非有界自己 Ext と同じく
$\mathcal P\notin\mathcal C_A$ である。
また $\operatorname{RHom}(\mathcal O,\mathcal P)=\mathbb C$ なので、
$\mathcal P$ 自体は残余圏の対象ではない。
\end{remark}

\subsection{二ブロックの定常点を残余圏へ射影する}

次の Serre 入力は全 coherent 圏の smoothness を要求しない。

\begin{lemma}[混合型の線束圏上の Serre 双対]
\label{cmo:lem:mixed-line-serre}
$N=3k+1$、$n=N-2$、$a=N-3$ とする。
本節の混合中心三次商の proper 圏 $\mathcal C_A$ は
\[
 S_{\mathcal C_A}=(-a)[n]
\]
を Serre 関手として持つ。特に線束 perfect な定常点では
$S_{\mathcal C_A}(\mathcal P)\simeq\mathcal P[n]$ である。
\end{lemma}
\begin{proof}
命題\ref{cmo:prop:proper-contraction}の証明は、
各因子の potential 解消の tensor と中心差 $u_L-u_R$ の Koszul 解消から
\[
 R\underline{\operatorname{Hom}}_{B\otimes B^{\mathrm{op}}}
     (B,B\otimes B^{\mathrm{op}})\simeq B(N)[-N]
\]
を真正な次数付き双加群として構成している。
この構成は中心商の係数 $c_0$ に依存しない。
正則中心三次による change of rings と両側局所双対により
\[
 R\Gamma_A(A)\simeq A^\vee(a)[-(N-1)]
\]
を得る。
これは、定理\paperref{moduli}{thm:nfm-relative-serre}の証明で使う
局所双対式\papereqref{moduli}{eq:nfm-relative-local-dual}と同じ入力であり、
一ブロックという表示はこの入力を得た後には使わない。
同証明の導来局所化の三角を双対化すると、tails の双対化双加群は
内部 shift $-a$ と cohomological shift $N-2$ を持つ。
従って全 compact 対象 $K,L$ について
\[
 R\operatorname{Hom}(K,L)^\vee
 \simeq R\operatorname{Hom}(L,K(-a)[N-2])
\]
となる。有限中心模型では compact 部分は $\mathcal C_A$ であり、
この圏は proper であるから、これが主張した Serre 双対である。
$u:P(-1)\to P$ の余核は有限次元なので
$\mathcal P(-1)\simeq\mathcal P$ であり、最後の結論を得る。
\end{proof}

\begin{proposition}[残余射影の自己 Ext]
\label{cmo:prop:projected-point-ext}
$k=2$ とし、少なくとも一つの $\kappa(p_i)$ が非零である定常点
$\mathcal P$ を取る。
$\mathcal E=\langle\mathcal O,\mathcal O(1),\mathcal O(2),\mathcal O(3)\rangle$ と置き、
\[
 F\longrightarrow\mathcal P\longrightarrow K
\]
を $F\in\mathcal E$、$K\in\mathcal E^\perp\cap\mathcal C_A
=\mathcal K_A^{\mathrm{line}}$ となる射影三角とする。
このとき
\begin{equation}\label{cmo:eq:projected-point-ext}
 \operatorname{Ext}^j(K,K)=
 \begin{cases}
  \mathbb C,&j=0,\\
  \mathbb C^{10},&j=1,\\
  0,&j\ne0,1.
 \end{cases}
\end{equation}
特に $\chi(K,K)=-9$ であり、$K$ は exceptional 対象でも spherical 対象でもない。
この結論は全六十四個の定常点に適用できる。
\end{proposition}
\begin{proof}
\textbf{第一段階：斉次分解と例外射影。}
一因子では $X=x-au,Y=y-bu,Z=z-cu$ と置き、
式\eqref{cmo:eq:character-resolution}の行列の定数 $a,b,c$ を
$au,bu,cu$ へ置換する。これにより
\[
 0\to\mathcal B_\alpha(-3)\to\mathcal B_\alpha(-2)^3
 \to\mathcal B_\alpha(-1)^3\to\mathcal B_\alpha
 \to\mathbb C[u]\to0
\]
を得る。$u=0$ では $S(1,-1,-1)$ の Koszul 解消となる。
増大複体の全項が $\mathbb C[u]$ 上平坦であり、内部次数は下に有界なので、
$u$ による帰納法で完全性が持ち上がる。
二因子の $\mathbb C[u]$ 上の tensor は、右加群の規約へ移せば
\[
 B(-i)^{\binom6i}\qquad(0\leq i\leq6)
\]
を第 $i$ 項とする $P$ の有限自由 $B$-分解となる。

斉次な higher homotopy を使う Shamash 分解 $Q_\bullet\to P$ の項は
\[
 Q_m=\bigoplus_{q\ge0}A(-m-q)^{\binom6{m-2q}},
\]
ただし範囲外の二項係数は零とする。
$Q_\bullet(3)$ のうち、自由生成元の line twist が非負の項を取ると、
微分で閉じた有限部分複体
\begin{equation}\label{cmo:eq:point-exceptional-complex}
 F=\left[
 \mathcal O^{20}\longrightarrow
 \mathcal O(1)^{15}\oplus\mathcal O\longrightarrow
 \mathcal O(2)^6\longrightarrow\mathcal O(3)
 \right]
\end{equation}
を得る。表示した項の cohomological 次数は順に $-3,-2,-1,0$ である。
増大写像と $\mathcal P(3)\simeq\mathcal P$ により $F\to\mathcal P$ がある。

$0\leq i,j\leq3$ では線束間の高次 Ext は零である。
従って $R\operatorname{Hom}(\mathcal O(i),F)$ は項ごとの Hom で計算できる。
これは自由分解 $Q_\bullet$ の内部次数 $3-i$ の成分そのものである。
実際、$F$ から除いた項の line twist は負なので、この成分に寄与しない。
従ってこの Hom 複体は次数零の $P_{3-i}=\mathbb C$ のみを持つ。
また exact な $j^*\dashv j_*$ により
\[
 R\operatorname{Hom}(\mathcal O(i),\mathcal P)
 =R\operatorname{Hom}_D(D,M)=\mathbb C
\]
であり、上の増大写像はこの同型を与える。
よって全ての $i=0,1,2,3$ について $R\operatorname{Hom}(\mathcal O(i),K)=0$。
命題\ref{cmo:prop:point-perfect}により $\mathcal P$ は線束 perfect なので、
$K$ も線束 perfect である。
従って式\eqref{cmo:eq:point-exceptional-complex}は所要の例外射影を実際に表す。

\textbf{第二段階：$R\operatorname{Hom}(F,\mathcal P)$ と制限写像。}
二因子の affine ambient character 分解に Hom を施した複体を
$(C^\bullet,\delta)$ と書く。
各因子は次元 $(1,3,3,1)$、唯一の非零微分が二つの三次元空間の間の
$\operatorname{diag}(-2a_i,-2b_i,-2c_i)$ である。
従って tensor 複体では
\[
 \dim(C^0,C^1,C^2,C^3)=(1,6,15,20),\qquad
 \delta:C^1\hookrightarrow C^2,\qquad
 \operatorname{rank}(\delta:C^2\to C^3)=9.
\]
全 Shamash cochain 複体は
$C^\bullet[t]$、$|t|=2$、微分
$\delta+t h_1^*+t^2h_2^*+\cdots$ である。
式\eqref{cmo:eq:point-exceptional-complex}に対応する Hom 複体は
\[
 C^0\longrightarrow C^1\longrightarrow C^2\oplus tC^0
 \longrightarrow C^3.
\]
最初の微分は零、次の微分は $\delta:C^1\to C^2$ の成分が単射なので階数六、
最後の微分は $\delta:C^2\to C^3$ なので階数九である。
従って
\begin{equation}\label{cmo:eq:exceptional-to-point-ext}
 \operatorname{Ext}^j(F,\mathcal P)=
 \begin{cases}
 \mathbb C,&j=0,2,\\
 \mathbb C^{11},&j=3,\\
 0,&\text{その他}.
 \end{cases}
\end{equation}

自然な写像
\[
 \operatorname{Ext}^j(\mathcal P,\mathcal P)
 \longrightarrow\operatorname{Ext}^j(F,\mathcal P)
\]
は、全 Shamash cochain 複体から上の有限複体への制限で計算できる。
これは $u\ne0$ への制限と $j^*\dashv j_*$ の随伴により確認できる。
次数零では恒等写像である。
次数二では両複体の次数一・二の項が同じなので、cohomology 上の写像は単射である。
両側の次元は一だから同型となる。
次数三の単射性も直接確認できる。
制限が境界となる全 cocycle から、その境界の持上げを引くと、
残りは $t v$、$v\in C^1$ の形である。
その微分の $tC^2$ 成分は $t\delta v$ である。
$\delta:C^1\to C^2$ が単射だから $v=0$ となる。
従って次数三でも制限写像は単射である。

\textbf{第三段階：proper 線束圏での双対と cone。}
前補題を $k=2$ に適用すると $S=(-4)[5]$ であり、
$\mathcal P(-4)\simeq\mathcal P$ である。
従って式\eqref{cmo:eq:exceptional-to-point-ext}を双対化して
\[
 \operatorname{Ext}^j(\mathcal P,F)=
 \begin{cases}
 \mathbb C^{11},&j=2,\\
 \mathbb C,&j=3,5,\\
 0,&\text{その他}
 \end{cases}
\]
を得る。Serre 双対の自然性により
$\operatorname{Ext}^j(\mathcal P,F)\to
\operatorname{Ext}^j(\mathcal P,\mathcal P)$ は、
第二段階の次数 $5-j$ の写像の双対である。
従って次数二では全射、次数三・五では同型となる。
この双対を適用している圏は proper な $\mathcal C_A$ であり、
全 $D^b(\operatorname{qgr}A)$ の proper 性や smoothness は使わない。

$R\operatorname{Hom}(F,K)=0$ なので射影三角から
\[
 R\operatorname{Hom}(K,K)\simeq R\operatorname{Hom}(\mathcal P,K)
 \simeq\operatorname{Cone}\bigl(
 R\operatorname{Hom}(\mathcal P,F)\to
 R\operatorname{Hom}(\mathcal P,\mathcal P)\bigr).
\]
既存の自己 Ext の次元 $(1,1,1,1)$（次数 $0,2,3,5$）と、
今得た写像の階数を長完全列へ代入すると、
次数零に一、次数一に $11-1=10$ のみが残る。
これが式\eqref{cmo:eq:projected-point-ext}である。
\end{proof}

\begin{remark}[数値的な独立検算と積について]
例外列の Euler 行列は
\[
 U=\begin{pmatrix}
 1&7&28&83\\0&1&7&28\\0&0&1&7\\0&0&0&1
 \end{pmatrix}.
\]
$\chi(\mathcal O(i),\mathcal P)=1$ から
$U(c_0,c_1,c_2,c_3)^{\mathsf t}=(1,1,1,1)^{\mathsf t}$ を解くと
$(c_0,c_1,c_2,c_3)=(-19,15,-6,1)$ となり、
\[
 [K]=[\mathcal P]+19[\mathcal O]-15[\mathcal O(1)]
                 +6[\mathcal O(2)]-[\mathcal O(3)]
\]
を得る。$\chi(\mathcal P,\mathcal P)=0$、
$\chi(\mathcal P,\mathcal O(i))=-1$ と合わせると $\chi(K,K)=-9$ であり、
上の全次数計算と一致する。
この証明は、自己 Ext の計算に使った外積表示を Yoneda 積と同定していない。
また $\mathcal P$ の自己 Ext と残余射影 $K$ の自己 Ext は異なる。
\end{remark}


\subsection{定常 norm とは異なる六十四個の非滑らかな係数}

\begin{theorem}[Azumaya 部分の Morse 点から生じる中心商の特異性]
\label{cmo:thm:morse-singular-sixty-four}
例\ref{cmo:ex:sixty-four}の二因子を取る。
命題\ref{cmo:prop:center-critical-polynomials}で定める四次多項式を $Q_1(W),Q_2(W)$ とし、
\begin{equation}\label{cmo:eq:morse-resultant}
 P_{64}(c)=\operatorname{pp}_+
   \operatorname{Res}_{z}\bigl(Q_1((-c-100z)^2),Q_2(z^2)\bigr)
       \in\mathbb Z[c]
\end{equation}
と置く。$\operatorname{pp}_+$ は整数係数の最大公約数を除き、
最高係数を正とする規格化である。
この多項式は次数六十四の平方因子を持たない偶多項式であり、
\[
 P_{64}(0)\ne0,\qquad \gcd(P_{64},N)=1
\]
を満たす。ここで $N$ は定常 character の norm である。

$P_{64}(c)=0$ の各係数に対して、$A_c$ の中心 chart は
三十六次元単純右加群 $M$ を持ち、その完備中央近傍は
\begin{equation}\label{cmo:eq:morse-matrix-node}
 \operatorname{Mat}_{36}\!\left(
   \C[[v_1,\ldots,v_6]]/(v_1^2+\cdots+v_6^2)\right)
\end{equation}
に同型である。その coherent な延長 $\mathcal P(M)=\pi(M[u])$ は
\begin{equation}\label{cmo:eq:morse-self-ext}
 \sum_{j\ge0}\dim_\C\operatorname{Ext}_{\operatorname{qgr}A_c}^j
       (\mathcal P(M),\mathcal P(M))s^j
       =\frac{(1+s)^6}{1-s^2}
\end{equation}
を満たす。
従って、これら全六十四係数で qgr の homological dimension は無限、
$D^b(\operatorname{qgr}A_c)$ は非 proper、線束生成圏とその残余部分は非 smooth であり、全 coherent 残余圏も非 proper である。
特に、この六十四係数は滑らかな可換 cubic への全 coherent Ku 比較が
成立する係数には含まれない。
\end{theorem}
\begin{proof}
$Q_i$ の根は零でなく相異なり、対応する $\Omega_i$ の八個の値
\[
 \mathcal M_i=\{\lambda\in\C:Q_i(\lambda^2)=0\}
\]
は、中心の正則部分における非退化臨界点の値である
（命題\ref{cmo:prop:center-critical-polynomials}）。
終結式の定義より、$P_{64}$ の根集合は
$-\mathcal M_1-100\mathcal M_2$ である。
次数、平方因子の不存在、零および norm 根との不交性は、
二つの四次多項式から有理数体上の終結式と最大公約数を計算して確認する。
再現コードは \texttt{verify\_mixed\_center\_geometry.py} であり、
係数の浮動小数点近似や根の数値分離は使わない。
次数六十四で平方因子がないため、八行八列の和の六十四値は全て相異なる。

$c=-\lambda-100\mu$ に対応する二つの臨界点 $z_1,z_2$ を取る。
定理\ref{cmo:thm:azumaya-regular}から、各 $D_{\alpha^{(i)}}$ は
$z_i$ の近傍で次数六の Azumaya 代数である。
その外積は、滑らかな六次元中心の $z=(z_1,z_2)$ の近傍で
次数三十六の Azumaya 代数になる。
$\Omega_1+100\Omega_2+c$ の Hessian は二つの非退化な三次元 Hessian の
block 対角和であり、係数 $100\ne0$ なので非退化である。
標数零の形式的 Morse 補題により、その方程式は形式座標変更の下で
$v_1^2+\cdots+v_6^2$ となる。
Azumaya 代数を剰余体 $\C$ の完備局所環上で分裂させ、中心商を取れば
式\eqref{cmo:eq:morse-matrix-node}を得る。

この行列代数の単純剰余加群は次元三十六である。
完備化前にも、対応する中心点上の唯一の単純加群 $M$ として定義される。
局所有限自由分解と完備化の平坦性により、$M$ の自己 Ext は完備化後と一致する。
実際、中心の極大イデアルは Hom 複体の標的 $M$ を零化し、
全ての自己 Ext も同じ剰余体上のベクトル空間なので、完備化で変わらない。

Morita 同値の下でこれは五次元通常二重点の剰余体の自己 Ext である。
六変数正則局所環の Koszul 分解に、極大イデアルの二乗に属する一つの二次式の
Shamash 変数を加えた分解は極小であり、その Betti 級数は
$(1+s)^6/(1-s^2)$ となる。
補題\ref{cmo:lem:chart-ext}で全次数の qgr 自己 Ext へ移すと
式\eqref{cmo:eq:morse-self-ext}を得る。

\paragraph{有限中央台に対する両方向の Ext の局所性.}
補題\ref{cmo:lem:finite-support-perfect}の有限中心模型
$Z=\operatorname{Proj}C^{(\ell)}$、$\mathcal A=\widetilde{A^{(\ell)}}$ を使う。
有限次元 chart 加群 $M$ の延長 $\mathcal P=\mathcal P(M)$ の中央台は、
$U=D_+(u^\ell)$ に含まれる有限閉集合である。
従って任意の coherent $\mathcal A$-加群 $\mathcal F$ について
\begin{equation}\label{cmo:eq:finite-support-local-ext}
 \operatorname{Ext}^j_{\operatorname{qgr}A}(\mathcal P,\mathcal F)
 \simeq\operatorname{Ext}^j_D(M,j^*\mathcal F)
 \qquad(j\ge0)
\end{equation}
となる。実際、各層 $\mathcal E xt^j_{\mathcal A}(\mathcal P,\mathcal F)$ は
coherent で、その台はこの有限閉集合に含まれる。
したがって高次層 cohomology は零であり、局所-to-大域 Ext spectral sequence は
次数ごとに退化する。その大域切断は $U$ 上の切断に等しく、
$U$ は affine なので右辺となる。
この議論は $\mathcal P$ の perfect 性を仮定しない。
有限中心模型と coherent 部分の導来完全忠実性は既存の補題と同じものを使っている。

\paragraph{節点上の単純加群から coherent 残余圏の非 proper 性へ.}
式\eqref{cmo:eq:morse-matrix-node}の単純加群 $M$ に対し、
\[
 R=\mathbb C[[v_1,\ldots,v_6]]/(v_1^2+\cdots+v_6^2)
\]
と置く。六変数正則局所環の Koszul 双対と中心正則元の Rees 同型から
\[
 R\operatorname{Hom}_R(\mathbb C,R)\simeq\mathbb C[-5]
\]
である。$\operatorname{Mat}_{36}(R)$ の右正則加群は Morita 同値で
$R^{36}$ に対応し、その単純加群は $\mathbb C$ に対応する。
したがって局所化・完備化の忠実平坦性により
\[
 \operatorname{Ext}^j_D(M,D)=
 \begin{cases}
  \mathbb C^{36},&j=5,\\
  0,&j\ne5.
 \end{cases}
\]
有限生成加群の有限階数自由分解を用いれば、各次数の Ext は完備化と可換である。
またその中央台は $M$ の一点の台に含まれるため、
この一点での局所化・完備化の計算が元の Ext を決める。
式\eqref{cmo:eq:finite-support-local-ext}と $j^*\mathcal O(i)\simeq D$ から、
全ての $i$ について
\begin{equation}\label{cmo:eq:node-point-to-line}
 \operatorname{Ext}^j(\mathcal P,\mathcal O(i))=
 \begin{cases}
  \mathbb C^{36},&j=5,\\
  0,&j\ne5
 \end{cases}
\end{equation}
を得る。
逆方向は exact な chart の随伴から
$R\operatorname{Hom}(\mathcal O(i),\mathcal P)\simeq M$ である。

有限例外列
$\mathcal E=\langle\mathcal O,\mathcal O(1),\mathcal O(2),\mathcal O(3)\rangle$
に沿って有限 mutation を行い、
\[
 F\longrightarrow\mathcal P\longrightarrow K,
 \qquad F\in\mathcal E,\quad
 K\in\mathcal K_A^{\mathrm{coh}}=\mathcal E^\perp
\]
を取る。線束から coherent 対象への $R\operatorname{Hom}$ は有限なので、
この三角は $D^b(\operatorname{qgr}A)$ の中で構成される。
式\eqref{cmo:eq:node-point-to-line}と有限 cone・直和因子に対する閉性により
$R\operatorname{Hom}(\mathcal P,F)$ は有界で有限次元である。
一方 $R\operatorname{Hom}(F,K)=0$ だから、射影三角より
\begin{equation}\label{cmo:eq:node-residual-self-ext}
 R\operatorname{Hom}(K,K)
 \simeq
 \operatorname{Cone}\bigl(
  R\operatorname{Hom}(\mathcal P,F)\longrightarrow
  R\operatorname{Hom}(\mathcal P,\mathcal P)\bigr).
\end{equation}
従って十分高い全ての次数で
\[
 \operatorname{Ext}^j(K,K)
 \simeq\operatorname{Ext}^j(\mathcal P,\mathcal P)
 \simeq\mathbb C^{32}.
\]
最後の次元は
$(1+s)^6/(1-s^2)=(1+s)^5/(1-s)$ の係数から従う。
このように全 coherent 残余圏自体が非 proper である。
また線束圏は proper なので、この $K$ は残余線束圏に属さない。
滑らかな可換 cubic の Ku は proper なので、両者は同値にならない。
この結論は、全 $D^b(\operatorname{qgr}A)$ の dg smoothness 単独を否定する主張ではない。

\paragraph{線束圏とその残余部分について.}
式\eqref{cmo:eq:morse-self-ext}から
$D^b(\operatorname{qgr}A)$ が非 proper であることは直ちに従う。
線束生成圏 $\mathcal C_A$ は常に proper であり、
もし smooth なら定理\paperref{main}{crs:thm:coherent}により全 coherent 圏へ飽和して矛盾する。
従って $\mathcal C_A$ は非 smooth である。
また命題\paperref{main}{crs:prop:residual-smooth}による有限例外 gluing から、
残余線束圏が smooth なら $\mathcal C_A$ も smooth になるので、
残余線束圏も非 smooth である。

\end{proof}

\begin{remark}[二種類の六十四係数]
定常 norm の六十四根は有理数であり、そこで生じる一次元定常点は
命題\ref{cmo:prop:point-perfect}により全て線束 perfect である。
これに対し $P_{64}$ の六十四根では、三十六次元単純加群が
非有界自己 Ext を持つ。この二集合は交わらない。
従って定常点の出現値だけでは非 smooth locus を捉えられないことが、
同じ固定混合族の中で具体的に示された。
定理\ref{cmo:thm:exact-sixty-four}で節点切断と射影境界を調べると、
非 smooth 係数は正確に $P_{64}$ の六十四根である。
系\ref{cmo:cor:fixed-zero-norm-smooth}により、
norm の六十四係数と $c=0$ は全て smooth proper である。
それらでの可換 Ku 比較は定理\ref{cmo:thm:all-central-commutative-kernel}と
系\ref{cmo:cor:norm-commutative-kernel}で構成され、主予想への反例にはならない。
\end{remark}
